%!TEX root = ../main.tex

% ============================================================
%  Kapitel 6: Testing und Evaluation
%  ARis -- AR-Brille Studienarbeit | Bastian Kiefer
% ============================================================

\chapter{Testing und Evaluation}
\label{chap:TestingundEvaluation}

Die Validierung des entwickelten Systems erfolgte auf Basis der in
Abschnitt~\ref{sec:Anforderungsanalyse} definierten Anforderungen.
Anstelle automatisierter Testsuites kamen aufgrund des prototypischen Charakters
des Projekts qualitative Funktionstests zum Einsatz, die sowohl begleitend während
der Entwicklung als auch in einer abschließenden Integrationsphase durchgeführt wurden.
Die Ergebnisse geben damit Auskunft darüber, in welchem Umfang die formulierten
Must-Have- und Should-Have-Anforderungen erfüllt wurden.

% ---------------------------------------------------------------------------
\section{Testmethodik}
\label{sec:testmethodik}

Die Testmethodik orientiert sich am prototypischen Charakter des Projekts:
Auf automatisierte Unit- oder Integrationstests wurde bewusst verzichtet, da der
Fokus auf einer möglichst breiten funktionalen Abdeckung des Gesamtsystems lag und
der verfügbare Zeitrahmen eine vollständige Testautomatisierung nicht zuließ.
Stattdessen wurden drei Testkategorien definiert, die zeitlich gestaffelt durchgeführt
wurden.

\emph{Begleitende Tests} fanden während der gesamten Entwicklungsphase statt:
Jede neu implementierte Komponente wurde
unmittelbar nach ihrer Fertigstellung funktional überprüft, bevor die Entwicklung
fortgesetzt wurde.
\emph{Komponentenspezifische Tests} wurden am Ende der jeweiligen Entwicklungsabschnitte
durchgeführt und prüften die Einzelkomponenten in isolation -- das Backend und das Frontend jeweils für sich.
\emph{Systemtests} erfolgten nach der Systemintegration und überprüften das
Zusammenspiel aller Komponenten im Gesamtbetrieb auf der Zielhardware.

Die Bewertung der Testergebnisse erfolgt dreistufig:
\textbf{Erfolgreich} bezeichnet eine vollständige und anforderungsgemäße Funktion,
\textbf{Teilweise erfolgreich} bezeichnet eine grundsätzlich funktionierende, jedoch
in Umfang oder Qualität eingeschränkte Funktion, und
\textbf{Nicht umgesetzt} bezeichnet Funktionen, die im Projektzeitraum nicht
implementiert oder getestet werden konnten.

% ---------------------------------------------------------------------------
\section{Optische Tests}
\label{sec:optischeTests}

Die optischen Tests wurden begleitend während des mechanischen Zusammenbaus des
Gesamtsystems durchgeführt.
Da optische Parameter wie Brennweite, Abstand und Spiegelwinkel im Prototypenbau
nicht vollständig rechnerisch vorausberechnet werden konnten, wurden diese iterativ
am realen System erprobt und angepasst.
Tabelle~\ref{apx:optischetests} fasst die durchgeführten Tests und ihre Ergebnisse
zusammen.

Die Testergebnisse zeigen, dass die grundlegende optische Funktion des
Free-Space-Combiner-Systems realisiert werden konnte: Das Display leuchtet,
das Bild wird durch das Okular kollimiert, durch den Umlenkspiegel geführt
und auf dem Combiner sichtbar gemacht, während die Umgebung gleichzeitig
durchsichtig bleibt.
Die wesentlichen Einschränkungen liegen in der Bildgröße und der Fokussierung,
die auf die Kombination aus dem gewählten Okular und dem mechanischen Aufbau
zurückzuführen sind.

% ---------------------------------------------------------------------------
\section{Software-Tests}
\label{sec:SoftwareTests}

Die Software-Tests gliederten sich in begleitende Entwicklungstests -- d.\,h.\
manuelle Prüfungen nach jeder abgeschlossenen Implementierungseinheit -- sowie
einen abschließenden manuellen Gesamttest auf der Zielhardware (Raspberry~Pi~Zero~2\,W
mit angeschlossenem Display).
Tabelle~\ref{apx:softwaretests} dokumentiert die Ergebnisse der wesentlichen
Testbereiche.

% ---------------------------------------------------------------------------
\section{Systemtest}
\label{sec:SystemTest}

Nach Abschluss der Integrationsphase wurde ein qualitativer Gesamtsystemtest
auf der physischen Zielplattform durchgeführt.
Dabei wurden alle implementierten Softwaremodule gemeinsam im realen Betrieb
auf dem Raspberry~Pi~Zero~2\,W erprobt.
\newpage
Der Test bestätigte, dass das Gesamtsystem als integrierte Einheit stabil läuft:
Backend und Frontend starten gemeinsam, der Datenpfad von der Companion App über
das Backend bis zur Darstellung auf dem Display funktioniert, und die
Taster-Eingabe wird korrekt verarbeitet.
Trotz der zuvor beschriebenen Deployment-Schwierigkeiten
(vgl.~Abschnitt~\ref{sec:softwareImplementierung}) -- insbesondere dem erforderlichen
OS-Downgrade und den damit verbundenen Paketkompatibilitätsproblemen -- läuft das
System nach der Einrichtung stabil und reproduzierbar.

Die optische Integration zeigte die in Abschnitt~\ref{sec:optischeTests} beschriebenen
Einschränkungen: Das angezeigte Bild ist kleiner als konzeptionell vorgesehen, und
die Fokussierung ist nicht für simultane Nah- und Fernwahrnehmung optimiert.
Die Demo-Seite konnte jedoch als funktionierender Proof of Concept auf der
physischen Hardware demonstriert werden -- ein AR-ähnliches HUD ist im Combiner-Spiegel
sichtbar und interaktiv durch den Taster bedienbar.

% ---------------------------------------------------------------------------
\section{Anforderungsabgleich}
\label{sec:anforderungsabgleich}

Abschließend werden die Testergebnisse den zu Beginn definierten MoSCoW-Anforderungen
gegenübergestellt.
Tabelle~\ref{tab:anforderungsabgleich} gibt einen Überblick über den
Erfüllungsgrad der Must-Have- und Should-Have-Anforderungen.

\begin{table}[htbp]
  \centering
  \caption{Abgleich der Testergebnisse mit den MoSCoW-Anforderungen
           (vgl.~Anhang~\ref{apx:anforderungen}).}
  \label{tab:anforderungsabgleich}
  \begin{tabularx}{\textwidth}{@{} l l X @{}}
    \toprule
    \textbf{Kategorie} & \textbf{Anforderung} & \textbf{Erfüllungsgrad} \\
    \midrule
    Must Have & Optische Sichtbarkeit der Anzeige &
      Grundsätzlich erfüllt; Bildgröße eingeschränkt \\
    Must Have & Freihand-Nutzbarkeit &
      Erfüllt -- kein Handgerät für Betrieb nötig \\
    Must Have & Mobile Verwendbarkeit &
      Erfüllt -- Akku-Betrieb realisiert \\
    Must Have & Eingabemöglichkeit (Taster) &
      Erfüllt \\
    Must Have & Verschiedene Darstellungsmodi &
      Überwiegend erfüllt (Hub, Demo, Teleprompter, Messages) \\
    Must Have & Physische Stabilität &
      Erfüllt -- Gehäuse hält Betrieb stand \\
    Must Have & Kosteneffizienz &
      Erfüllt -- Realisierung mit handelsüblichen Komponenten und innerhalb des Budgets \\
    \midrule
    Should Have & Brillenähnliches Design &
      Erfüllt \\
    Should Have & Gyroskop / Lageerfassung &
      Erfüllt -- MPU-6050 integriert und auslesbar \\
    Should Have & Temperatur-/Feuchtigkeitssensoren &
      Überwiegend Erfüllt -- GY-63 integriert und auslesbar erweiterbar durch z.B. BME280 \\
    Should Have & Companion App &
      Überwiegend erfüllt (WiFi-Kanal funktionsfähig) \\
    Should Have & Möglichst hohe Displayauflösung &
      Bedingt erfüllt -- Bildbereich eingeschränkt \\
    \bottomrule
  \end{tabularx}
\end{table}

Der Anforderungsabgleich zeigt, dass alle Must-Have-Anforderungen zumindest
grundsätzlich erfüllt wurden -- mit der Einschränkung, dass die optische Bildgröße
hinter den konzeptionellen Erwartungen zurückbleibt.
Sämtliche Should-Have-Anforderungen wurden ebenfalls überwiegend realisiert.
Die nicht umgesetzten Features (Feature Apps, Pixel Streaming, vollständige
BLE-Verbindung) entstammen ausschließlich der Could-Have-Kategorie oder waren
als konzeptionelle Vorbereitungen angelegt, ohne im Projektzeitraum vollständig
realisiert werden zu können.